\chapter{A model artifact-evaluation section}
\label{app:aec-model}

Chapter~\ref{ch:repro} argues that artifact evaluation is where an intrinsic identifier earns
its keep: a reviewer can check an artifact offline, in under a minute, without asking anyone's
permission. This appendix is the operational form of that argument --- text a programme chair
can lift into a call for artifacts, and two short procedures, one for authors and one for
reviewers.

It is written to be \emph{additive}. Every venue that runs artifact evaluation already has a
deposit rule, a badge vocabulary, and a reviewing workload it is trying not to increase. Nothing
below removes a DOI, displaces Zenodo, or asks a committee to retrain its reviewers. It fills a
slot that most guidelines already contain and few can currently fill.

It states \emph{what} must be provided and \emph{what} is checked --- not the incantations.
Tools drift faster than principles, so the commands, the special cases and the troubleshooting
live in a companion \textsc{howto}, maintained and versioned alongside these notes
(\texttt{docs/HOWTO-artifact-swhid.md} in their repository,
\url{https://github.com/rdicosmo/source-code-of-science/blob/main/docs/HOWTO-artifact-swhid.md}).
The \textsc{howto} is deliberately referenced by path rather than pinned to one revision: it is
expected to move as tooling does, which is the entire point of keeping it out of this appendix.

\section{The slot that is already there}

The ACM Artifact Review and Badging policy asks, for the \emph{Artifacts Available} badge, that
authors provide \enquote{a DOI or link to this repository along with a unique identifier for the
object}. Two slots. The first is routinely filled with a DOI; the second is usually left as a bare
repository URL --- which is a location, not an identifier of an object. Many venues inherit this
sentence verbatim, and a good number have quietly dropped its second half.

\takeaway{Most guidelines already ask for an identifier of the object. Few can name one.}

Filling that second slot with a \textsc{swhid} costs an author two local commands and gives a
reviewer something they have never had: a claim about the artifact that can be checked rather
than trusted. Everything else in this appendix follows from that one substitution.

\section{What the identifier names}
\label{sec:aec-invariant}

One invariant carries the whole workflow. The identifier always names a \emph{clean tree}: for a
git-managed artifact, the committed tree of a precise commit; for anything else, the content of
the exact bundle submitted, unpacked into an empty directory. Nobody --- author, reviewer or
committee --- ever hashes a live working copy: a checkout contains \texttt{.git/}, build output
and editor droppings, so its hash matches neither the archive nor the reviewer's tarball, and is
not even stable over time on an unchanged commit. For a git artifact the clean tree's
\texttt{dir} \textsc{swhid} is, bit for bit, the hash that \verb|git rev-parse <commit>^{tree}|
already prints --- which is also exactly what Software Heritage stores.

Hence one recipe, current at press time (maintained version, variants and troubleshooting in the
\textsc{howto}; the local tools in Appendix~\ref{app:cookbook}):

\begin{quote}\footnotesize
\begin{verbatim}
d=$(mktemp -d) && git archive <commit> | tar -x -C "$d"
swh identify --no-filename --type directory "$d"
git rev-parse <commit>^{tree}   # self-check: the two hashes must be equal
\end{verbatim}
\end{quote}

\noindent If the self-check fails, the tree is not clean: diagnose (the \textsc{howto} has the
checklist) rather than submit the number.

One rule then orders the identifier's two possible sources. \emph{Before} archival --- at
submission and through revisions, double-blind included --- the clean-export value is the only
one there is: the artifact is not in the archive yet, and under double-blind it must not be.
\emph{After} archival, the qualified \textsc{swhid} from the archive's \emph{Permalinks} tab is
primary and the local value the cross-check; where both exist they must coincide, and that
coincidence is itself the end-to-end check of the whole procedure.

\takeaway{Everyone hashes the same object: the clean tree, never the working copy.}

\section{Model text for a call for artifacts}

\begin{callout}{Artifact identification and archival --- text to adapt}
\small
\textbf{Identifying the artifact.} At every submission or revision, authors must include the
\texttt{dir} \textsc{swhid} (\emph{SoftWare Hash IDentifier}, ISO/IEC~18670) of the \emph{clean
artifact tree}: the committed tree for a version-controlled artifact, the content of the
submitted bundle otherwise --- never a live working copy. The identifier is \emph{intrinsic}:
derived from the content alone, it requires no account, no deposit and no network access to
produce, and it changes whenever the artifact changes. Compute it following the \textsc{howto}
accompanying this call, which includes a one-line self-check. A revised artifact gets a new
\textsc{swhid}, reported with the revision; earlier identifiers remain valid names for the
earlier states.

\smallskip
\textbf{Archiving the artifact --- after acceptance.} In addition to any deposit this venue
already requires (for example Zenodo, with its DOI), authors must ensure the source code is
archived in Software Heritage: request archival of the repository --- history included ---
through Save Code Now (\url{https://save.softwareheritage.org}); no account is needed. Then take
the \emph{qualified} \textsc{swhid} of the exact version from the archive's \emph{Permalinks}
tab and put it in the camera-ready paper alongside the DOI of any deposit. Its core \texttt{dir}
hash must equal the identifier recorded at the close of evaluation; verify that equality before
submitting. If the artifact was improved after evaluation, list both identifiers, labelled
\enquote{as evaluated} and \enquote{as published}.

\smallskip
\textbf{Before acceptance there is nothing to archive.} Under double-blind review, do not
request archival before acceptance: the archive records the repository's origin, which would
reveal the authors. A core directory \textsc{swhid} carries no author, origin or commit
metadata, so it may appear in an anonymized submission as is; qualifiers (\texttt{origin},
\texttt{anchor}) are added at camera-ready around the identical core hash.

\smallskip
\textbf{What reviewers check.} Reviewers verify the identification by \emph{recomputation}: they
unpack the artifact into a fresh directory, recompute the identifier per the \textsc{howto}, and
compare it with the one on record. Equality proves that the tree under review is, bit for bit,
the tree the authors named. A mismatch goes through the \textsc{howto}'s diagnostic checklist
first --- its commonest cause is procedural, not a changed artifact --- and is reported as a
finding only if it survives diagnosis.

\smallskip
\textbf{Scope.} This covers \textbf{source code}, including build recipes --- Dockerfiles,
lockfiles, provisioning and experiment scripts, pinned under the same identifier as the artifact
tree. Datasets, container \emph{images} and large binaries follow this venue's existing
instructions; Software Heritage is a source-code archive, not a data archive.
\end{callout}

\section{For authors}

Two phases, and the phase determines where the identifier comes from.

\paragraph{Before archival --- submission and revisions.}
\begin{enumerate}[nosep,leftmargin=*]
  \item Develop wherever you already develop. Nothing changes about your forge, your CI or your
        packaging.
  \item At submission, compute the identifier of the clean artifact tree --- the recipe in
        Section~\ref{sec:aec-invariant}, or the \textsc{howto} for anything it does not cover ---
        and put it in the artifact appendix beside your submission URL. If the self-check fails,
        diagnose before submitting the number.
  \item At each revision during evaluation, recompute and report it in one line:
        \enquote{revised tree: \texttt{swh:1:dir:\dots}}.
\end{enumerate}

\paragraph{After acceptance.}
\begin{enumerate}[nosep,leftmargin=*]
  \item Archive the repository through Save Code Now, and take the qualified \textsc{swhid} of
        the exact release from the \emph{Permalinks} tab.
  \item Check that its core \texttt{dir} hash equals the identifier recorded at the close of
        evaluation. Equality here is the end-to-end check that every step above went right; on a
        mismatch, go to the \textsc{howto}'s diagnostics, or label the two versions
        \enquote{as evaluated} and \enquote{as published}.
  \item Deposit as your venue requires --- unchanged --- and cite both the DOI and the
        \textsc{swhid} in the camera-ready paper.
\end{enumerate}

\section{For reviewers}

What a submitted tarball supports is precise, and worth stating honestly.

\paragraph{What you can verify, offline.} That the tree you unpacked \emph{is} the object the
authors named: unpack into a fresh directory, recompute the identifier per the \textsc{howto},
compare. Equality is a proof of bit-for-bit identity, and the comparison involves no judgement,
no account, no service, and nothing beyond the initial download.

\paragraph{What a mismatch means.} Only that two different trees were hashed --- nothing about
\emph{how} different, or on whose side the error lies; the commonest cause is procedural (a live
working copy hashed on one side), not a changed artifact. Run the \textsc{howto}'s diagnostics
first; report only what survives them, as you would a build failure: the identifier obtained and
the procedure used, never an accusation.

\paragraph{What you cannot verify.} Before acceptance, anything about the archive: the artifact
is not there, and under double-blind it must not be --- the check is purely local. Nor that the
submitted tree corresponds to any commit in the authors' (hidden) repository; that link is
established only at camera-ready. For the \emph{Available} badge after acceptance, the real
check is to resolve the camera-ready qualified \textsc{swhid} --- resolution proves archival ---
and to confirm that its core \texttt{dir} hash equals the identifier recorded at the close of
evaluation. (Fully independent recomputation from the archive: see the \textsc{howto}.)

\takeaway{The comparison is arithmetic: recompute, compare. Interpreting a mismatch is not:
diagnose before you report.}

\section{How it maps onto existing badges}

\begin{center}\small
\begin{tabular}{@{}>{\raggedright\arraybackslash}p{3.2cm}>{\raggedright\arraybackslash}p{9.6cm}@{}}
\toprule
\textbf{Badge} & \textbf{What changes} \\
\midrule
Artifacts Available & The identifier slot is filled with a verifiable identifier. The DOI keeps
its role: citation metadata, discovery, a landing page. \\
\addlinespace
Artifacts Evaluated & The committee can state \emph{which} artifact it evaluated, and the
published version can be shown to be that one. \\
\addlinespace
Results Reproduced & Unaffected in substance, but the reproduction now names the exact tree it
reproduced from. \\
\bottomrule
\end{tabular}
\end{center}

\section{The questions a chair will ask}

\paragraph{We already require a DOI. Why a second identifier?}
They answer different questions. A DOI names a deposit and certifies a landing page; the deposit
behind it can be replaced without the DOI changing, and nothing in the identifier tells you it
did. A \textsc{swhid} \emph{is} the content. That is why both belong in the paper, not one
instead of the other.

\paragraph{Is Software Heritage durable enough to rely on?}
The identifier does not depend on the archive's survival: it can be recomputed from any copy of
the content, by anyone, indefinitely. This is the practical difference between an intrinsic
identifier and a registry entry --- the registry can fail; arithmetic cannot.

\paragraph{Does this cover datasets, containers and models?}
No. It identifies source code, and build recipes pinned with it. Datasets and images keep the
arrangements your venue already has. A proposal that claimed otherwise would be easy to refute
and would deserve to be.

\paragraph{What does it cost the committee?}
One sentence in the call for artifacts, one line in the reviewing form, and a link to the
\textsc{howto}. No new infrastructure, no accounts, no service dependency, and a verification
step that is faster than the one it supplements.
